% Principe d'un spectrophotomètre (L'essentiel, chapitre 2)
\begin{tikzpicture}[x=1cm,y=1cm,
  boite/.style={draw=otInk!70,fill=otRule!35,line width=0.6pt,rounded corners=1.2mm,minimum height=11mm,
    align=center,font=\otLight\fontsize{7.5}{9}\selectfont,text=otInk},
  faisceau/.style={line width=2.2pt,-{Stealth[length=2.6mm,width=4mm]}}]
\node[boite,minimum width=16mm] (s) at (0,0) {source\\de lumière};
\node[boite,minimum width=20mm] (m) at (2.9,0) {monochromateur\\(choix de $\lambda$)};
\pic at (5.9,-0.75) {cuve={hauteur=1.5,niveau=1.15,liquide=pcBleu!70}};
\node[boite,minimum width=16mm] (d) at (8.9,0) {détecteur};
\pic (af) at (11.3,-0.55) {afficheur={largeur=1.6,hauteur=1.1,texte={0,62}}};
\draw[faisceau,otOrange!80] (s.east) -- (m.west);
\draw[faisceau,otGreen] (m.east) -- (5.3,0) node[midway,above,pcLegende]{$I_0$};
\draw[faisceau,otGreen!55] (6.5,0) -- (d.west) node[midway,above,pcLegende]{$I<I_0$};
\draw[pcFil] (d.east) -- (10.5,0);
\draw[otMuted,line width=0.4pt,{Stealth[length=1.4mm]}-{Stealth[length=1.4mm]}] (5.46,-0.95) -- (6.34,-0.95)
  node[midway,below,pcLegende]{$\ell$};
\node[pcLegende,anchor=north] at (5.9,-1.25) {cuve contenant la solution};
\node[pcLegende,anchor=north] at (11.3,-0.65) {absorbance $A$};
\end{tikzpicture}
\par\smallskip\legende{La lumière monochromatique est en partie absorbée par la solution : $A=\log(I_0/I)$.}
